\documentclass[11pt]{article}

%%METADATA
\title{GA1025 Macro Theory \\Solution to Problem Set 8}
\author{
Junbiao Chen\thanks{E-mail: jc14076@nyu.edu.}
}
\date{\today}


%%PACKAGES
\usepackage{mdframed} % For boxed environments
\usepackage{graphicx}
\usepackage{grffile}
\usepackage{tabularx}
\usepackage{setspace}
\usepackage{amsmath,amsthm,amssymb}
\usepackage[hyphens]{url}
\usepackage{natbib}
\usepackage[font=normalsize,labelfont=bf]{caption}
\usepackage[margin=1in]{geometry}
\usepackage{hyperref}
\hypersetup{colorlinks=true,urlcolor=blue,citecolor=blue}
\usepackage{stmaryrd}  %Package with \boxast command
\usepackage{enumerate}% http://ctan.org/pkg/enumerate %Supports lowercase Roman-letter enumeration
\usepackage{verbatim} %Package with \begin{comment} environment
%\usepackage{enumitem}
\usepackage{physics}
\usepackage{tikz}
\usepackage{listings}
\usepackage{upquote}
\usepackage{booktabs} %Package with \toprule and \bottomrule
\usepackage{etoc}     %Package with \localtableofcontents
\usepackage{placeins}    %Package that prevent repositioning the tables
\usepackage{multicol}
\usepackage{bm}
\usepackage{subfig}
\usepackage{csquotes}

\definecolor{dkgreen}{rgb}{0,0.6,0}
\definecolor{gray}{rgb}{0.5,0.5,0.5}
\definecolor{mauve}{rgb}{0.58,0,0.82}

\lstset{language=bash,
  frame=tb,
  aboveskip=3mm,
  belowskip=3mm,
  showstringspaces=false,
  columns=flexible,
  basicstyle={\small\ttfamily},
  numbers=none,
  numberstyle=\tiny\color{gray},
  keywordstyle=\color{blue},
  commentstyle=\color{dkgreen},
  stringstyle=\color{mauve},
  breaklines=true,
  breakatwhitespace=false,
  tabsize=3
}

\lstset{language=C,
  aboveskip=3mm,
  belowskip=3mm,
  showstringspaces=false,
  columns=flexible,
  basicstyle={\small\ttfamily},
  numbers=none,
  numberstyle=\tiny\color{gray},
  keywordstyle=\color{blue},
  commentstyle=\color{dkgreen},
  stringstyle=\color{mauve},
  breaklines=true,
  breakatwhitespace=false,
  tabsize=4
}

\definecolor{lightblue}{rgb}{0.68, 0.85, 0.9} %

%CUSTOM DEFINITIONS
\theoremstyle{definition}
\newtheorem{definition}{Definition}[section]
\newtheorem*{remark}{Remark}
\setcounter{secnumdepth}{3}
\usepackage{tikz}
\usetikzlibrary{arrows.meta}
\usetikzlibrary{automata, positioning, arrows, calc}

\tikzset{
	->,  % makes the edges directed
	>=stealth, % makes the arrow heads bold
	shorten >=2pt, shorten <=2pt, % shorten the arrow
	node distance=3cm, % specifies the minimum distance between two nodes. Change if n
	every state/.style={draw=blue!55,very thick,fill=blue!20}, % sets the properties for each ’state’ n
	initial text=$ $, % sets the text that appears on the start arrow
 }

%% PROPOSITION
% Define the Proposition environment
\newmdenv[
  innerleftmargin=10pt, 
  innerrightmargin=10pt,
  innertopmargin=10pt,
  innerbottommargin=10pt,
  linecolor=black, 
  linewidth=1pt,
  backgroundcolor=white, 
  roundcorner=5pt
]{propositionbox}

\newtheoremstyle{boldtitle} % Define a new theorem style
  {10pt} % Space above
  {10pt} % Space below
  {\itshape} % Body font
  {} % Indent amount
  {\bfseries} % Theorem head font
  {.} % Punctuation after theorem head
  { } % Space after theorem head
  {} % Theorem head spec

\theoremstyle{boldtitle} % Use the custom style
\newtheorem{proposition}{Proposition} % Define the proposition environment

% Redefine the proposition environment to use the box
\newenvironment{boxedproposition}[1][]
{\begin{propositionbox}\begin{proposition}[#1]}
{\end{proposition}\end{propositionbox}}

%%FORMATTING
\usepackage[bottom]{footmisc}
\onehalfspacing
\numberwithin{equation}{section}
\numberwithin{figure}{section}
\numberwithin{table}{section}
\bibliographystyle{../bib/aeanobold-oxford}


%main text
\begin{document}

\maketitle


\section{Exercise 9.3 Existnece of representative consumer}
The Lagrangian is given by 
\[
\mathcal{L} = \theta u(c_1) + (1 - \theta) w(c_2) + \lambda (c - c_1 - c_2)
\]
The FOCs are 
\begin{align*}
    [c_1]: \quad \theta u'(c_1) & = \lambda \\
    [c_2]: \quad (1 - \theta) w'(c_2) & = \lambda
\end{align*}
It follows that 
\begin{equation}
    \label{eqn:foc_consumption}
    \theta u'(c_1) = (1 - \theta) w'(c_2)
\end{equation}

Since the both $u$ and $w$ are strictly increasing, at the optimal solution, we have 
\[
c_1(c, \theta) + c_2(c, \theta) = c,
\]
where $c_1$ and $c_2$ are optimal solutions. 

Taking derivative w.r.t. to $c$, we have 
\begin{equation}
    \label{eqn:foc_constraint}
c_{1,c}'(c, \theta) + c_{2,c}'(c, \theta) = 1
\end{equation}

Now plugging $c_1(c, \theta)$ and $c_2(c, \theta)$ into the objective function, we have 
\[
\theta u(c_1(c, \theta)) + (1- \theta) w(c_2(c, \theta)) = v_\theta (c)
\]
Taking derivative w.r.t. to $c$, we have 
\[
\theta u'(c_1(c, \theta)) c_{1,c}'(c, \theta) + (1- \theta) w'(c_2(c, \theta)) c_{2,c}'(c, \theta) = v_\theta' (c)
\]
Applying equations~(\ref{eqn:foc_consumption}) and (\ref{eqn:foc_constraint}), we have 
\[
(1- \theta) w'(c_2(c, \theta)) c_{1,c}'(c, \theta) + (1- \theta) w'(c_2(c, \theta)) (1 - c_{1,c}'(c, \theta)) = v_\theta' (c)
\]
and
\[
\theta u'(c_1(c, \theta)) (1 - c_{2,c}'(c, \theta)) + \theta u'(c_1(c, \theta)) c_{2,c}'(c, \theta) = v_\theta' (c)
\]
Simplifying these equations, we have 
\[
v_\theta' (c) = \theta u'(c_1(c, \theta)) =  (1- \theta) w'(c_2(c, \theta)),
\]
which represents the marginal utility of the representative consumer.



\section{Exercise 9.4 Term structure of interest rates}
Denote the state of the economy at time $t$ by $s_t \in \{s_1, s_2 \}$.
Then, the consumer's problem is given by%
\footnote{
  Since there is only one consumer, the superscript $i$ is not needed.
}
\begin{align*}
& \max_{c_t(s^t)} \quad \sum_{t=0}^{\infty} \sum_{s^t} \beta^t \pi(s^t) \frac{c_t(s^t)^{1 - \gamma}}{1 - \gamma} \\
& \text{subject to } \quad  \sum_{t=0}^{\infty} \sum_{s^t} q^0_t(s^t) c_t (s^t) = \sum_{t=0}^{\infty} \sum_{s^t} q^0_t(s^t) y_t (s^t)
\end{align*}
where $q^0_t(s^t)$ is the price of the time $t$ history $s^t$ goods at time 0.

\subsection{a. Definition of competitive equlibrium}
A competitive equilibrium is an allocation and a price system such that 
\begin{enumerate}
  \item Given the price system $q_t^0(s^t)$, the allocation solves the consumer's problem, and
  \item Market clears: $c_t(s^t) = y_t(s^t)$.
\end{enumerate}

\subsection{b. solution}
To compute the equilibrium is to derive the an allocation $c_t(s^t)$ and prices $q^0_t(s^t)$.
The Lagrangian for the consumer's problem is given by
\[
\mathcal{L} = \sum_{t=0}^{\infty} \sum_{s^t} \beta^t \pi(s^t) \frac{c_t(s^t)^{1 - \gamma}}{1 - \gamma} 
+ \mu \left( \sum_{t=0}^{\infty} \sum_{s^t} q_t^0(s^t) (y_t(s^t) - c_t(s^t)) \right)
\]
FOC is given by 
\[
\beta^t \pi(s^t) c_t(s^t)^{- \gamma} = \mu q_t^0(s^t) 
\]
To remove the Lagrangian multiplier, notice that $\pi(s^0) = 1$ (because the state at time 0 is given)
and standardize $q_0^0(s^0) =1$.
Then, we have 
\[
\beta^t \pi(s^t) c_t(s^t)^{- \gamma} = c_0^{- \gamma} q_t^0(s^t) 
\]
Applying the market clearing conditions, then we have 
\[
\beta^t \pi(s^t) y_t(s^t)^{- \gamma} = y_0^{- \gamma} q_t^0(s^t) 
\]
Since $y_t(s^t) = \prod_{j=1}^t \lambda_j(s_j) y_0$, we have 
\[
\beta^t \pi(s^t) \left( \prod_{j=1}^t \lambda_j(s_j) y_0 \right)^{- \gamma} = y_0^{- \gamma} q_t^0(s^t) 
\]
Therefore, in the competitive equilibrium, the price system is given by 
\[
q_t^0(s^t) = \beta^t \pi(s^t) \left( \prod_{j=1}^t \lambda_j(s_j) \right)^{- \gamma},
\]
and the allocation is given by 
\[
c_t(s^t) = c_0 \prod_{j=1}^t \lambda_j(s_j).
\]

\noindent \textbf{Remark:} This model is beautiful as prices and allocations arise endogenously and are of closed-formed.





%% =========== %% ========= %%
\bibliography{../bib/references.bib}

\end{document}